% Vantagens_participar_SPE.tex
\documentclass[11pt,a4paper]{article}
\usepackage[utf8]{inputenc}
\usepackage[T1]{fontenc}
\usepackage[portuguese]{babel}
\usepackage{lmodern}
\usepackage{geometry}
\geometry{margin=2.5cm,headheight=15pt}
\usepackage{microtype}
% Vantagens_participar_SPE.tex — versão limpa e formal
\documentclass[11pt,a4paper]{article}
\usepackage[utf8]{inputenc}
\usepackage[T1]{fontenc}
\usepackage[portuguese]{babel}
\usepackage{geometry}
\usepackage{hyperref}
\usepackage{enumitem}
\geometry{margin=2.5cm}

	itle{Documento de Apresentação — SPE Student Chapter ISPTEC}
\author{}
\date{Março de 2026 — Versão provisória}

\begin{document}
\maketitle

\section*{Estatuto do documento}
Esta nota explicativa foi preparada pela equipa promotora do Student Chapter SPE — ISPTEC. Trata-se de uma versão provisória; a circulação institucional só deverá ocorrer após validação técnica, revisão linguística e aprovação institucional.

\section{Resumo executivo}
O presente documento sintetiza os principais argumentos e benefícios da criação do Student Chapter SPE no ISPTEC e esclarece o estatuto atual do núcleo e do Board para fins de comunicação com stakeholders internos.

\section{Estatuto do Núcleo e do Board}
Em termos institucionais, o documento clarifica que o Núcleo se encontra em fase de instalação e que a inscrição formal está prevista para o período indicado pelos promotores. O Board atual tem caráter provisório até à confirmação formal pela instituição e pela SPE regional/internacional.

\section{Formato documental recomendado}
Recomenda-se que este texto seja assumido como \emph{Nota Explicativa} ou \emph{Memorando Académico} na sua versão para circulação institucional, com um cabeçalho que identifique claramente a versão (interna/provisória vs. oficial).

\section{Seleção de membros — critérios (resumo)}
Para transparência, sumariza-se aqui os critérios objectivos que devem suportar a selecção (ver relatório técnico para o quadro completo):
\begin{itemize}[leftmargin=1.0cm]
  \item Interesse demonstrado e motivação;
  \item Disponibilidade para participar nas actividades programadas;
  \item Situação da filiação SPE (membership) quando aplicável;
  \item Equilíbrio entre cursos e anos académicos;
  \item Critérios específicos para composição do Board (competências técnicas, experiência organizacional, diversidade de cursos);
  \item Prioridade estratégica a perfis com competências críticas nas fases iniciais (ex.: organização, contactos institucionais).
\end{itemize}

\section{Proteção de dados}
Antes da partilha institucional, remover contactos telefónicos e anonimizar identificadores (SPE ID, matrículas) sempre que não sejam estritamente necessários. Manter uma versão interna restrita com dados completos e uma versão limpa para circulação externa.

\section{Revisão e validação}
Antes de oficializar, o documento deve passar por:
\begin{enumerate}[leftmargin=1.0cm]
  \item Revisão técnica (conteúdo e factos); 
  \item Revisão linguística (ortografia e estilo); 
  \item Validação de conteúdos sensíveis (dados pessoais e afirmações institucionais);
  \item Alinhamento com os actores académicos directamente envolvidos (coordenações, Chapter Advisor).
\end{enumerate}

\section*{Próximos passos}
\begin{itemize}[leftmargin=1.0cm]
  \item Preparar a versão limpa para circulação institucional (redacção final e anonimização);
  \item Submeter para revisão técnica e linguística; 
  \item Agendar reunião de validação com Chapter Advisor e coordenações relevantes.
\end{itemize}

\vspace{1em}
\noindent Contato para coordenação: Chapter Advisor / equipa promotora (inserir contacto institucional no ficheiro interno).

\end{document}
\begin{itemize}
